\section{CART high-cardinality selection bias (idealized)}
\label{app:cart-bias}

The following symmetry calculation formalizes the multiple-comparisons
mechanism behind ``high-cardinality'' selection bias in greedy CART-style
split selection \citep{Breiman1984ClassificationAndRegressionTrees}. It is
\emph{motivational} and is not used for any of the fixed-node p-value
guarantees stated here.

\begin{lemma}[Threshold search amplifies null extremes (union bound)]
\label{lem:cart-union-bound}
Fix a node and a feature $j$ with a finite candidate threshold set $C_j$ of size
$m_j := |C_j| \ge 1$. For each $c \in C_j$, let $\Delta_{j,c}$ be a split score
for candidate split $(j,c)$, where larger values are better. Then for any
$\delta \in \mathbb{R}$,
%
\begin{equation}
  \PP\!\left(\max_{c \in C_j} \Delta_{j,c} \ge \delta\right)
  \le \sum_{c \in C_j} \PP(\Delta_{j,c} \ge \delta)
  \le m_j \cdot \max_{c \in C_j} \PP(\Delta_{j,c} \ge \delta).
\end{equation}
%
If additionally $(\Delta_{j,c})_{c \in C_j}$ are i.i.d.\ and satisfy
$\PP(\Delta_{j,c} \ge \delta) = p_\delta$, then
%
\begin{equation}
  \PP\!\left(\max_{c \in C_j} \Delta_{j,c} \ge \delta\right) = 1-(1-p_\delta)^{m_j}.
\end{equation}
\end{lemma}

\begin{proof}
The first inequality is the union bound:
$\PP(\max_{c \in C_j}\Delta_{j,c} \ge \delta)
= \PP(\cup_{c \in C_j}\{\Delta_{j,c} \ge \delta\})
\le \sum_{c \in C_j}\PP(\Delta_{j,c} \ge \delta)$.
The second inequality is immediate.
If the scores are i.i.d.\ with tail probability $p_\delta$, then
$\PP(\max_{c \in C_j}\Delta_{j,c} < \delta) = \prod_{c \in C_j} (1-p_\delta)
  = (1-p_\delta)^{m_j}$, so taking complements yields the identity.
\end{proof}

\begin{remark}[Categorical split counts can be exponential]
If a greedy method treats an unordered categorical feature with $L$ levels as
allowing all binary partitions of the levels, the number of distinct candidate
splits is $2^{L-1}-1$ (non-empty proper subsets, modulo swapping left/right).
Many implementations avoid enumerating all such partitions, but this
combinatorial explosion highlights why ``high-cardinality'' can create very
large multiple-comparisons effects in greedy split optimization.
\end{remark}

\begin{lemma}[Proportional selection under an exchangeable null]
\label{lem:cart-proportional-selection}
Assume A0.6 (\cref{app:assumptions}). Consider a fixed node with candidate
features $j \in \{1,\dots,p\}$ and candidate split sets $C_j$ with sizes
$m_j := |C_j| \ge 1$. Let $S_{j,c}$ denote a split score for candidate split
$(j,c)$, where smaller is better, and assume the collection
$\{S_{j,c} : 1 \le j \le p,\ c \in C_j\}$ is exchangeable.

Let $(V_{j,c})$ be i.i.d.\ $\mathrm{Unif}(0,1)$ independent of the scores and
define tie-broken scores $\widetilde S_{j,c} := (S_{j,c}, V_{j,c})$ ordered
lexicographically. Define the greedy CART choice
%
\begin{equation}
  (\widehat j, \widehat c) \in \argmin_{1 \le j \le p,\ c \in C_j} \widetilde S_{j,c}.
\end{equation}
%
Then for each feature $j$,
%
\begin{equation}
  \PP(\widehat j = j) = \frac{m_j}{\sum_{k=1}^p m_k}.
\end{equation}
\end{lemma}

\begin{proof}
By construction, the augmented collection $\{(S_{j,c},V_{j,c})\}$ is
exchangeable and almost surely tie-free. Therefore the argmin is uniformly
distributed over the $\sum_{k=1}^p m_k$ candidate splits, and it belongs to
feature $j$ with probability $m_j/\sum_{k=1}^p m_k$.
\end{proof}

\begin{remark}[Feature subsampling does not remove the mechanism]
If a greedy CART-style method first samples a subset of candidate features
$F \subseteq \{1,\dots,p\}$ (e.g., as in random forests) and then minimizes the
score over candidates within $F$, the same symmetry argument gives, conditional
on $F$,
%
\begin{equation}
  \PP(\widehat j = j \mid F) = \frac{m_j}{\sum_{k \in F} m_k} \qquad (j \in F).
\end{equation}
%
Thus high-cardinality features remain favored \emph{within} the tested subset.
\end{remark}
